\begin{lemma}[\textup{(5.8)}: the right-hand side term]
Let $n\ge1$, let $p>1$ and let $\gamma_4>0$.  Let $\Omega\subseteq\R^n$, let
$(A,B)$ satisfy (Q3) of the thesis--(Q5) of the thesis on $\Omega$ with the parameters
$p,\varepsilon,\mu,\gamma_2,\gamma_3,\gamma_4$
(\noderef{EllipticStructureQ345}), let $u,v$ be real valued and $g$ a vector
field, let $x_0\in\R^n$ and $R>0$ with $B_R(x_0)\subseteq\Omega$, let
$\alpha>0$ and $C_1\ge0$, and assume
\[
g\in L^p\bigl(B_R(x_0)\bigr),
\qquad
B(\cdot,u,g)\in L^1\bigl(B_R(x_0)\bigr),
\qquad
|u-v|\le4C_1R^{\alpha}\ \text{ a.e. on }B_R(x_0),
\]
and assume further that $u$ and $v$ are almost everywhere strongly measurable
on $B_R(x_0)$.
Then
\[
\Bigl|\int_{B_R(x_0)}B\bigl(x,u(x),g(x)\bigr)\bigl(u(x)-v(x)\bigr)dx\Bigr|
\ \le\ 4C_1\gamma_4R^{\alpha}
\int_{B_R(x_0)}\bigl(1+|g|^2\bigr)^{\frac p2}dx .
\]
\end{lemma}
\begin{proof}
Write $B_R=B_R(x_0)$ and note $|B_R|<\infty$ (Preamble: the volume of a ball).

\emph{Step 1: the growth majorant is integrable.}  By clause \texttt{Q5} of
\noderef{EllipticStructureQ345}, for almost every $x\in\Omega$, every $z$ and
every $h$,
\[
\bigl|B(x,z,h)\bigr|\le\gamma_4\bigl(1+|h|^2\bigr)^{\frac p2} .
\]
Since $B_R\subseteq\Omega$ this holds for almost every $x\in B_R$; read at
$z=u(x)$ and $h=g(x)$,
\[
\bigl|B(x,u(x),g(x))\bigr|\le\gamma_4\bigl(1+|g(x)|^2\bigr)^{\frac p2}
\qquad\text{a.e. on }B_R .
\tag{1}
\]
The function $(1+|g|^2)^{p/2}$ is integrable on $B_R$: by cases on $|g|\le1$,
where $1+|g|^2\le2$ so the power is at most $2^{p/2}$, and on $|g|>1$, where
$1+|g|^2\le2|g|^2$ so the power is at most $2^{p/2}|g|^p$ (pm-0024, the bases
being positive); hence
$(1+|g|^2)^{p/2}\le2^{p/2}(1+|g|^p)$, which is integrable on $B_R$ because
$g\in L^p(B_R)$ and $|B_R|<\infty$.  It is also nonnegative, so its integral is
a nonnegative real.

\emph{Step 2: the pointwise bound on the integrand.}  Combining the hypothesis
$|u-v|\le4C_1R^{\alpha}$ almost everywhere on $B_R$ with (1), and noting
$4C_1R^{\alpha}\ge0$ because $C_1\ge0$ and $R^{\alpha}\ge0$ (pm-0024),
\[
\bigl|B(x,u(x),g(x))\bigl(u(x)-v(x)\bigr)\bigr|
\le4C_1R^{\alpha}\,\gamma_4\bigl(1+|g(x)|^2\bigr)^{\frac p2}
\qquad\text{a.e. on }B_R .
\tag{2}
\]

\emph{Step 3: integrating.}  The integrand
$x\mapsto B(x,u(x),g(x))(u(x)-v(x))$ is integrable on $B_R$.  It is almost
everywhere strongly measurable on $B_R$, being the product of two such
functions: $B(\cdot,u,g)$, which is almost everywhere strongly measurable
because it lies in $L^1(B_R)$ by hypothesis, and $u-v$, which is almost
everywhere strongly measurable because $u$ and $v$ are, by hypothesis, and a
difference of two such functions is again one.  \emph{The measurability of $u$
and $v$ is a hypothesis and is not derivable from the rest of the list}: the
$L^1$ hypothesis constrains only the superposition $B(\cdot,u,g)$, the $L^p$
hypothesis only $g$, and the essential bound on $|u-v|$ carries no regularity
at all.  By (2) the integrand is dominated almost everywhere by
$4C_1R^{\alpha}\gamma_4(1+|g|^2)^{p/2}$, which is integrable on $B_R$ by Step
1; in particular that majorant is itself almost everywhere strongly measurable,
being a continuous function of $|g|$ and $g\in L^p(B_R)$ being almost
everywhere strongly measurable.  A function almost everywhere strongly
measurable and dominated almost everywhere by an integrable function is
integrable, so the integrand is integrable on $B_R$.  Hence, chaining
$|\int f|\le\int|f|$ with (2) and taking the constant out,
\[
\Bigl|\int_{B_R}B(x,u,g)(u-v)\,dx\Bigr|
\le\int_{B_R}\bigl|B(x,u,g)(u-v)\bigr|dx
\le4C_1\gamma_4R^{\alpha}\int_{B_R}\bigl(1+|g|^2\bigr)^{\frac p2}dx ,
\]
which is the assertion.

\emph{On the constant.}  The source's printed (5.8) of the thesis carries
$2C_1\gamma_4$, and its adjudicated \texttt{CHECK} at l.~1514 records why the
printed chain really needs $4C_1\gamma_4$: the hypothesis
$\sup|u(x_0)-u(x)|\le C_1R^{\alpha}$ gives
$\operatorname*{osc}_{B_R}u\le2C_1R^{\alpha}$, so the displayed chain for $|w|$
through a boundary point yields $4C_1R^{\alpha}$ rather than $2C_1R^{\alpha}$.
The present statement takes the essential bound on $|u-v|$ as a hypothesis at
the constant $4C_1R^{\alpha}$ and carries that constant into the conclusion, so
no step here asserts the printed inequality.  A consumer that can produce the
sharper essential bound $2C_1R^{\alpha}$ --- comparing $u$ and $v$ through the
single value $U(x_0)$ rather than through a boundary point --- may of course
weaken it to $4C_1R^{\alpha}$ and apply this lemma unchanged.
\end{proof}
